% Báo cáo tiến độ Đồ án 1 — biên dịch bằng XeLaTeX:
%   latexmk -xelatex bao_cao_tien_do.tex
\documentclass[11pt,a4paper]{article}

\usepackage[a4paper,top=2.2cm,bottom=2.2cm,left=2.5cm,right=2cm]{geometry}
\usepackage{fontspec}
\usepackage[vietnamese]{babel}
\setmainfont{texgyretermes}[
  Extension=.otf,
  UprightFont=*-regular,
  BoldFont=*-bold,
  ItalicFont=*-italic,
  BoldItalicFont=*-bolditalic]
\setmonofont{DejaVu Sans Mono}[Scale=0.85]

\usepackage{xcolor}
\usepackage{graphicx}
\usepackage{booktabs}
\usepackage{tabularx}
\usepackage{array}
\usepackage{enumitem}
\usepackage{titlesec}
\usepackage{fancyhdr}
\usepackage{tikz}
\usetikzlibrary{positioning,arrows.meta,fit,calc}
\usepackage[hidelinks]{hyperref}

% ---------------------------------------------------------------- thông tin
\newcommand{\tendetai}{HUST Course Advisor \& Degree Audit Engine}


\titleformat{\section}{\large\bfseries}{\thesection.}{0.5em}{}
\titlespacing*{\section}{0pt}{1.1em}{0.5em}
\titleformat{\subsection}{\normalsize\bfseries}{\thesubsection.}{0.5em}{}
\titlespacing*{\subsection}{0pt}{0.7em}{0.3em}

\setlist{itemsep=1pt,topsep=3pt,leftmargin=1.4em}
\setlength{\parindent}{0pt}
\setlength{\parskip}{0.45em}
\renewcommand{\arraystretch}{1.25}

\pagestyle{fancy}
\fancyhf{}
\fancyfoot[C]{\small\thepage}
\renewcommand{\headrulewidth}{0pt}

\newcommand{\xong}{Xong}
\newcommand{\dang}{Đang làm}
\newcommand{\code}[1]{\texttt{#1}}

\begin{document}

% ---------------------------------------------------------------- bìa
% Bìa theo mẫu tham_khao/report/src/frontmatter/cover.tex
\newgeometry{top=2cm,bottom=2cm,left=2cm,right=2cm}
\begin{titlepage}
\thispagestyle{empty}
\begin{tikzpicture}[remember picture,overlay]
  \draw[line width=1.0pt]
    ($(current page.north west)+(1.55cm,-1.35cm)$)
    rectangle
    ($(current page.south east)+(-1.55cm,1.35cm)$);
  \draw[line width=0.4pt]
    ($(current page.north west)+(1.63cm,-1.43cm)$)
    rectangle
    ($(current page.south east)+(-1.63cm,1.43cm)$);
\end{tikzpicture}
\begin{center}
  {\bfseries\fontsize{21}{26}\selectfont
  ĐẠI HỌC BÁCH KHOA HÀ NỘI\\
  TRƯỜNG ĐIỆN -- ĐIỆN TỬ\par}

  \vspace{0.15cm}
  \rule{6cm}{0.4pt}

  \vspace{0.8cm}
  \includegraphics[height=4.8cm]{logo.png}

  \vspace{0.7cm}
  {\bfseries\fontsize{21}{26}\selectfont BÁO CÁO TIẾN ĐỘ\\[0.25cm]
  ĐỒ ÁN 1\par}

  \vspace{1.0cm}
  {\bfseries\fontsize{18}{24}\selectfont\MakeUppercase{\tendetai}\par}

  \vspace{0.3cm}
  {\itshape\fontsize{13}{17}\selectfont Hệ thống chatbot tư vấn chọn môn, tra cứu quy chế\\
  và rà soát tiến độ học tập\par}

  \vfill

  {\fontsize{13}{18}\selectfont
  \begin{tabular}{@{}>{\bfseries}l@{\hspace{1.2em}}l@{}}
    Giảng viên hướng dẫn: & Bùi Tiến Cường \\
    Sinh viên thực hiện: & Mai Duy Hiếu \\
    MSSV: & 20223967 \\
  \end{tabular}\par}

  \vfill
  {\bfseries\fontsize{13}{17}\selectfont Hà Nội, 10/2026}
\end{center}
\end{titlepage}
\restoregeometry
\setcounter{page}{1}

% ---------------------------------------------------------------- 1
\section{Tóm tắt}

Đề tài kết hợp hai hướng thầy gợi ý: \emph{(i)} hỏi đáp quy chế đào tạo, học phí, học bổng
bằng RAG + Gemini và \emph{(ii)} chatbot tư vấn chọn môn theo CTĐT và điều kiện đăng ký.
Hiện tại hệ thống đã chạy được đầu--cuối (giao diện web $\rightarrow$ backend
$\rightarrow$ dữ liệu) cho ngành ET1, hai phiên bản CTĐT 2022 (K67) và
2023 (K68). Kho RAG gồm 5 văn bản chính thức đã đối chiếu với PDF gốc
(218 đoạn).

% ---------------------------------------------------------------- 2
\section{Kiến trúc đã triển khai}

\begin{center}
\begin{tikzpicture}[
  node distance=5mm and 9mm,
  box/.style={draw, rounded corners=2pt, align=center, font=\small,
              minimum height=8mm, inner sep=4pt, fill=white},
  db/.style={box, fill=black!6},
  arr/.style={-{Stealth[length=2mm]}, thick}]
  \node[box] (ui) {Frontend\\React 19 + Vite};
  \node[box, right=of ui] (api) {FastAPI\\Router phân loại ý định};
  \node[box, right=12mm of api, yshift=11mm] (rag) {RAG\\Retriever lai + Gemini};
  \node[box, right=12mm of api, yshift=-11mm] (rule) {Rule Engine\\+ Recommender};
  \node[db, right=of rag] (chroma) {ChromaDB\\5 văn bản, 218 đoạn};
  \node[db, right=of rule] (ctdt) {CTĐT JSON\\ET1-2022, ET1-2023};
  \node[db, below=of api] (sql) {SQLite\\lịch sử hội thoại};
  \draw[arr, <->] (ui) -- (api);
  \draw[arr] (api.east) -- ++(4mm,0) |- (rag.west);
  \draw[arr] (api.east) -- ++(4mm,0) |- (rule.west);
  \draw[arr] (rag) -- (chroma);
  \draw[arr] (rule) -- (ctdt);
  \draw[arr, <->] (api) -- (sql);
\end{tikzpicture}
\end{center}

% ---------------------------------------------------------------- 3
\section{Kết quả đã đạt được}

\begin{tabularx}{\linewidth}{@{}>{\raggedright\arraybackslash}p{3.1cm}X>{\centering\arraybackslash}p{1.6cm}@{}}
\toprule
\textbf{Hạng mục} & \textbf{Kết quả cụ thể} & \textbf{Trạng thái} \\
\midrule
Dữ liệu CTĐT &
  ET1-2022: 126 TC, 117 học phần, 20 nhóm, 6 mô-đun.
  ET1-2023: 132 TC, 79 học phần, 13 nhóm, 5 mô-đun.
  Điều kiện tiên quyết/học trước/song hành lấy từ CTT, giữ nguyên nhánh AND/OR. & \xong \\
Kho tài liệu RAG &
  5 văn bản đã kiểm chứng: CTĐT ET1 2022 và 2023; Quy chế đào tạo 2025
  (QĐ 5445/QĐ-ĐHBK); Sổ tay sinh viên 2026; QĐ học phí 2025--2026.
  Mỗi đoạn lưu metadata loại văn bản, thời gian hiệu lực, URL nguồn. & \xong \\
Pipeline RAG &
  Chia đoạn theo Điều/Khoản $\rightarrow$ ChromaDB $\rightarrow$ truy xuất lai
  (vector + từ khóa + khớp chính xác mã môn/Điều) $\rightarrow$ xếp hạng lại theo
  năm CTĐT và loại văn bản $\rightarrow$ Gemini 2.5 Flash. Khi không có API key
  vẫn trả lời bằng trích dẫn nguyên văn. & \xong \\
Rule Engine &
  Kiểm tra điều kiện đăng ký; rà soát tốt nghiệp (Degree Audit); tính CPA,
  hạn mức tín chỉ, mức cảnh báo; gợi ý môn kỳ tới theo mô-đun chuyên sâu. & \xong \\
Chatbot &
  Router phân loại ý định, hiểu tên gọi tắt của sinh viên (``vxl'', ``dts'', \dots).
  Lịch sử hội thoại lưu SQLite theo MSSV (ghim, lưu trữ, xóa). & \xong \\
Giao diện web &
  Đã có khung các tab Chat, Tiến độ (Audit), Gợi ý môn, Lộ trình; đang hoàn thiện giao diện. & \dang \\
Thu thập toàn trường &
  Danh mục 68 chương trình; rà link CTĐT (46 nguồn truy cập được); trích xuất
  14 bản nháp CTĐT và đối chiếu với 8.310 học phần trên CTT. & \dang \\
\bottomrule
\end{tabularx}

\subsection{Công việc trong thời gian gần đây}
\begin{itemize}
  \item Thay 3 bản quy chế tự soạn (chưa kiểm chứng) bằng văn bản trích từ PDF chính thức.
        Bản cũ được gắn \code{UNVERIFIED} và không còn được đưa vào chỉ mục.
  \item Cải thiện bộ truy xuất: thêm tìm theo từ khóa và xếp hạng lại, nên câu hỏi
        ``ET1 2023'' không còn bị trả về tài liệu 2022. Tự động xóa đoạn cũ khi văn bản thay đổi.
  \item Sửa router: câu dạng ``kỳ tới học X thì cần môn gì trước'' được chuyển sang kiểm tra
        điều kiện học phần thay vì gợi ý môn.
  \item Bỏ module đăng nhập tự viết (chưa đủ an toàn), dùng hồ sơ demo. Cấu hình địa chỉ
        API của frontend qua biến môi trường.
\end{itemize}

% ---------------------------------------------------------------- 4
\section{Khó khăn và hướng xử lý}

\begin{tabularx}{\linewidth}{@{}>{\raggedright\arraybackslash}p{5.2cm}X@{}}
\toprule
\textbf{Khó khăn} & \textbf{Cách đang xử lý} \\
\midrule
Dữ liệu điều kiện học phần trên CTT phải đăng nhập và chứa thông tin cá nhân. &
  Thu thập bằng script mở trình duyệt để tự đăng nhập. Chỉ lưu phần cấu trúc CTĐT,
  không công khai HTML gốc. \\
Nguồn CTĐT các ngành không đồng nhất (HTML, PDF, XLSX; nhiều link hỏng; chỉ 6/8 PDF
  đọc được danh mục học phần). &
  Xây pipeline tạo bản nháp và duyệt thủ công từng CTĐT trước khi đưa vào Rule Engine.
  Hiện mới duyệt xong ET1-2023. \\
Quy chế 2025 đổi thuật ngữ (``cảnh cáo'' $\rightarrow$ ``cảnh báo''; học bổng loại A/B/C),
  sinh viên vẫn hỏi theo cách cũ. &
  Thêm bảng từ đồng nghĩa khi truy vấn. \\
Tìm kiếm ngữ nghĩa với tiếng Việt hay nhầm Điều có số liệu gần nhau hoặc nhầm năm CTĐT. &
  Kết hợp tìm theo từ khóa và xếp hạng lại. Nhược điểm: phần xếp hạng hiện còn nhiều
  luật viết tay, khó mở rộng. \\
Gemini bản miễn phí bị giới hạn lượt gọi, model cũ bị ngừng hỗ trợ. &
  Cấu hình model qua \code{.env}. Khi không gọi được LLM, hệ thống chuyển sang trả
  trích dẫn nguyên văn có nguồn. \\
Chưa có cách đo chất lượng câu trả lời một cách định lượng. &
  Dự kiến xây bộ câu hỏi đánh giá. \\
\bottomrule
\end{tabularx}


\end{document}
